\chapter{次年度計画}

\figfull[.97\textwidth]{next_year_roadmap.pdf}{次年度の優先順位。まず評価を信頼できる形にし、基礎モデルと条件判断を強化し、強化学習の増分を検証した後、ミリメートル級の位置決めと接触制御が必要な管挿入洗浄へ進む。}

\begin{landscape}
\section{P0：結果の信頼性を確立する}
\begin{table}[H]
\centering\footnotesize
\renewcommand{\arraystretch}{1.35}
\caption{P0 作業パッケージ}
\begin{tabularx}{\linewidth}{>{\bfseries}p{24mm}p{48mm}p{65mm}p{55mm}X}
\toprule
目標 & 実施内容 & 検証指標 & 前提・リスク & 成果物\\
\midrule
試験手順の固定 & ボトル型、個数、姿勢、密度、照明、回収箱、反復回数を規定 & 各試験に場面ID、開始終了、逐次結果、異常理由 & 学習に入れない固定条件が必要；途中で判定規則を変えない & 標準試験手順書、固定試験一覧\\
成否判定 & 把持、キャップ把持、本体保持、回転、完全離脱、正しい箱への投入を別々に記録 & 工程別の自動・人手一致率、全工程成功率、95\%信頼区間 & 単独観測者の偏り；二名確認または映像確認 & 工程別結果表、自動集計画面\\
運用台帳 & 時刻、モデル版、データ版、ボトルID、結果、復旧を各周期で記録 & 欠測ゼロ、展示からモデルと条件へ追跡可能 & 個人・設備情報を秘匿 & 統一実験台帳、日次集計\\
入力取得から制御出力までの性能 & 推論遅延、制御欠落、動作列の切替、装置停止を記録 & 遅延p50/p95/p99、実制御周波数、切替時の指令値の不連続、停止回数 & 計測自体がリアルタイム性を損ねないこと & 性能基準、警報閾値\\
\bottomrule
\end{tabularx}
\end{table}

\vspace{3mm}
\section{P1：基礎モデルの成功率を高める}
\begin{table}[H]
\centering\footnotesize
\renewcommand{\arraystretch}{1.35}
\caption{P1 作業パッケージ}
\begin{tabularx}{\linewidth}{>{\bfseries}p{24mm}p{48mm}p{65mm}p{55mm}X}
\toprule
目標 & 実施内容 & 検証指標 & 前提・リスク & 成果物\\
\midrule
空回しの抑制 & 同型ボトルでキャップ有無以外の条件をそろえた対照軌跡を収録し、キャップなしは開栓省略を明示 & キャップなしで誤って開栓へ進む率、キャップあり見逃し率、全工程成功率 & 指示と実動作を一致させ、ボトル型の丸暗記を防ぐ & キャップ有無条件データセット、比較報告\\
逆向き把持の抑制 & ボトル口の向きと把持可能面で層別収集し、把持後の向き確認と再把持を追加 & 逆向き把持率、開栓可能率、再把持成功率 & フレーム単位の向き・把持端ラベルが必要 & 向き別の件数を均衡化したデータセット、把持診断表\\
接触安定化 & 接近、挟持、回転、離脱を分け、グリッパ電流と画像上の回転角を記録 & キャップ把持率、回転開始率、完全離脱率、落下率 & キャップ材質・摩擦差、画像角度の校正 & 接触工程データ、工程別比較\\
将来50制御周期分の動作列の改善 & 実行長、再計画時刻、異常時中断方式を比較 & 完了時間、切替時の指令値の不連続、復旧率、制御欠落 & データとモデルを同時変更しない & 動作列の条件比較報告\\
多視点頑健性 & 他条件を固定し、上部・手首視野の欠損または遮蔽を比較 & 3カメラ条件別の工程完了率、動作偏差 & 注意割合は遮蔽実験の代替ではない；安全確保 & カメラ寄与報告、最小センサ構成\\
\bottomrule
\end{tabularx}
\end{table}
\end{landscape}

\begin{landscape}
\section{P2：比較可能な強化学習へ進む}

次年度は基礎モデルと固定試験集合を先に凍結し、同じデータ、ボトル、初期条件、試験回数で基礎モデルと強化学習モデルを比較する。

\begin{table}[H]
\centering\footnotesize
\renewcommand{\arraystretch}{1.35}
\caption{P2 作業パッケージ}
\begin{tabularx}{\linewidth}{>{\bfseries}p{28mm}p{52mm}p{66mm}p{56mm}X}
\toprule
目標 & 実施内容 & 検証指標 & 前提・リスク & 成果物\\
\midrule
同一条件での基準モデル比較 & 基礎モデルと試験集合を固定し、強化学習側では一要因だけを変更 & 全工程・工程別成功率差、信頼区間、時間、接触異常 & 基礎モデルが継続変更されると比較不能 & 再確認可能な基準モデル\\
評価信号の改善 & 人手で付与した終端ラベルを工程別の評価ラベルへ分解し、重複・誤りを点検 & ラベル一致率、価値推定校正、オフライン順位正解率 & 誤った評価信号は誤動作を増幅 & 評価規約、洗浄済みデータ\\
複数の乱数初期値 & 重要設定を少なくとも3種類の初期値で反復 & 平均、標準偏差、95\%信頼区間 & 計算費用が高く、候補設定の事前絞込みが必要 & 反復条件の比較表\\
安全な実機確認 & オフライン選別後に小規模実機試験、動作境界と非常停止を設定 & 基礎モデルに対する工程別の改善幅、異常回数、復旧率 & 衝突と装置摩耗を最小化 & 強化学習実機評価報告\\
\bottomrule
\end{tabularx}
\end{table}
\end{landscape}

\begin{landscape}
\section{P3：空ボトルの管挿入洗浄}

次年度は空ボトルの口を水管へ挿入して洗浄する工程へ展開する。この課題は開栓よりも、口部中心、管中心、軸方向、接触時の追従性を厳密に要求する。「ミリメートル級」は目標であり、既に測定済みの精度ではない。

\begin{table}[H]
\centering\footnotesize
\renewcommand{\arraystretch}{1.35}
\caption{P3 作業パッケージ}
\begin{tabularx}{\linewidth}{>{\bfseries}p{28mm}p{52mm}p{66mm}p{56mm}X}
\toprule
目標 & 実施内容 & 検証指標 & 前提・リスク & 成果物\\
\midrule
実環境の幾何キャリブレーション & ボトル口、管口、カメラ、ロボット座標を測定し、反復位置決め誤差を記録 & 位置・角度誤差の平均、95パーセンタイル、再現性 & ボトル型の寸法差、校正ドリフト & 幾何基準、誤差予算\\
デジタルツイン & ボトル、水管、作業台、双腕の接触・衝突モデルを構築し、審査場面を固定 & 写真・測定との整合誤差、衝突一致率、軌跡実行可能率 & 外観一致は接触物理の一致ではない & 再確認可能な仮想試験環境\\
段階的挿入 & 位置決め、軸合わせ、接触、短距離挿入、洗浄完了に分解 & 工程別成功率、最大接触力、挿入深さ、完了時間 & 力覚がない場合は低速と視覚フィードバックを優先 & 段階挿入方式、安全規約\\
データと強化学習 & ミリメートル級の失敗境界と復旧を収録し、安全な仮想・オフライン環境で先に最適化 & 基礎模倣モデルに対する挿入誤差、復旧率、実機性能の改善幅 & 仮想・実機差、評価信号が危険動作を誘発する可能性 & 挿入データ、比較実験、論文資料\\
\bottomrule
\end{tabularx}
\end{table}
\end{landscape}

\section{次年度の主要な受入れ基準}

\begin{enumerate}
  \item すべての候補モデルを固定試験集合で全工程・工程別に評価できること。
  \item 空回しと逆向き把持に条件別誤り率を設定し、データまたは手法変更で有意に低減すること。
  \item 基礎模倣モデルと強化学習モデルを同条件・複数の乱数初期値で比較すること。
  \item 管挿入に実環境の幾何誤差予算、段階別指標、安全試験を設けること。
  \item すべての結果図を元データ、モデル版、試験台帳へ追跡できること。
\end{enumerate}
